%% ch02 --- Srodowiska i pakowanie: uv, pyproject i plik blokady
%%
%% Napisane. Kontraktem tego rozdzialu jest wpis ch02 w tools/chapters.json.
%% Kazdy listing to plik w code/ch02/, ktory uruchamia build, kazde cytowane
%% wyjscie to wygenerowany transkrypt, a kazda liczba trafia na strone przez
%% \val{} z code/measure/e2_install.py.

\chapter{Środowiska i pakowanie: uv, pyproject i plik blokady}
\label{ch:packaging}

\noindent
W .NET środowisko jest właściwością maszyny i nigdy nie musiałeś o nim
myśleć. SDK instaluje się raz, \code{dotnet restore} zapełnia jeden cache
pakietów, z którego czyta każdy projekt na maszynie, a wersja środowiska
uruchomieniowego, której chce projekt, to linia w pliku honorowana przez
narzędzia bez proszenia. Działa to tak niezawodnie, że jest niewidoczne.

Python odwraca to i na tym odwróceniu stoi cały rozdział.
\textbf{Środowisko Pythona to katalog.} Nie klucz w rejestrze, nie instalacja
ogólnosystemowa i nie ustawienie, które jakieś narzędzie odczytuje; to folder
z interpreterem w środku, a jedyne, co czyni go \emph{twoim} środowiskiem, to
fakt, że polecenie, które uruchomiłeś, akurat na niego wskazywało. Nic w
języku tego układu nie wymusza, nic nie ostrzega, gdy jesteś poza nim, a każdy
sposób, w jaki idzie to źle dla inżyniera .NET, wynika z tego jednego faktu.

\section{Mapa}
\label{sec:ch02-map}
\index{pakowanie!mapa narzędzi}

\noindent
Wszystko w tym rozdziale to jedna z tych par. Lewą kolumnę już znasz; prawa to
to, czego używa ta książka, a tam, gdzie jest wybór, powód wyboru znajdziesz w
kolejnej sekcji.

\begin{center}
\small
\begin{tabularx}{\linewidth}{@{}lX@{}}
\toprule
\textbf{.NET} & \textbf{Python, tak jak używa go ta książka} \\
\midrule
NuGet & PyPI, z \pkg{uv} jako klientem \\
\code{dotnet restore} & \code{uv sync} \\
jeden cache pakietów na maszynę & jeden katalog \code{.venv/} na projekt \\
\code{.csproj} & \code{pyproject.toml} \\
\code{packages.lock.json} & \code{uv.lock} \\
\code{dotnet run} & \code{uv run} \\
\code{dotnet tool} & \code{uvx} \\
analizatory plus \code{dotnet format} & \pkg{ruff} \code{check} i \pkg{ruff} \code{format} \\
\code{Program.cs} & punkt wejścia w \code{[project.scripts]} \\
jeden folder na projekt w rozwiązaniu & układ \code{src/} \\
\dash{} & \pkg{pre-commit}, który nie ma odpowiednika w .NET \\
\bottomrule
\end{tabularx}
\end{center}

\noindent
Jeden wiersz tej tabeli jest prawdziwą luką, a nie tłumaczeniem.
\pkg{pre-commit} uruchamia formatter i linter jako hak gita, na plikach, których
dotyka commit, zanim commit powstanie. .NET nie ma powszechnie używanego
odpowiednika, bo analizatory działają w buildzie, a build działa bez przerwy.
Build Pythona nie działa bez przerwy \dash{} nie ma buildu \dash{} więc to hak
musi być źródłem tej samej presji.

\section{Środowisko to katalog}
\label{sec:ch02-directory}
\index{środowisko!wirtualne}
\index{pakowanie!uv}

\noindent
Zanim przeczytasz dalej, odpowiedz na to ze swojego obecnego modelu Pythona.
Wpisujesz \code{pip install httpx} na maszynie, na której jest Python. Gdzie
wyląduje \pkg{httpx} i które z tego, co możesz potem uruchomić, będzie go
widzieć?

Uczciwa odpowiedź brzmi: nie da się tego wiedzieć bez sprawdzenia, i o to
właśnie chodzi. Nie ma jednego miejsca. \code{pip} instaluje do tego
interpretera, który go wywołał, a który to interpreter, zależy od twojego
\code{PATH}, od tego, czy w tej powłoce jest aktywne środowisko, i od tego, co
menedżer pakietów twojego systemu zrobił z \code{/usr/bin} lata temu.

Ta książka nigdy więc nie prosi, żebyś wiedział. Pyta interpreter:

\pyfile{code/ch02/where_am_i.py}{Środowisko opisane przez interpreter, który w nim siedzi}{lst:ch02-where}

\noindent
Uruchomione z \code{code/} wypisuje to:

\transcript{ch02-where-am-i}

\noindent
Cztery fakty, a każdy z nich jest pytaniem o katalog. \code{sys.prefix} to
korzeń środowiska; \code{sys.base\_prefix} to korzeń interpretera, z którego
środowisko zbudowano. Gdy się różnią, jesteś w środowisku wirtualnym, i to
porównanie jest całym testem \dash{} nie ma nic więcej do sprawdzenia.
Interpreter jest w środku katalogu. Pakiet zewnętrzny też jest w środku, w
\code{site-packages}. Biblioteki standardowej tam nie ma, bo jedzie razem z
interpreterem i nigdy nie była niczyją zależnością.

\begin{csbox}
\code{.venv/} jest najbliżej \code{obj/} i \code{bin/} razem wziętych: wynik
buildu, osobno na projekt, bezpieczny do skasowania, nigdy niecommitowany.
Różnica, która się liczy, jest taka, że skasowanie \code{obj/} kosztuje cię
przebudowę, a skasowanie \code{.venv/} nie kosztuje nic, bo \code{uv sync}
odbuduje je z pliku, który akurat zacommitowałeś. Traktuj je jak
jednorazowe, a za każdym razem będziesz miał rację.
\end{csbox}

\begin{trapbox}
\textbf{\code{pip install} kładzie pakiet na maszynie, jak narzędzie globalne.}
Nie kładzie, a wersja tego przekonania, która żyje najdłużej, jest tą groźną:
że nie ma wielkiego znaczenia, który interpreter go dostanie. Zainstaluj do
interpretera systemowego, a zmieniłeś zależność skryptów własnego systemu
operacyjnego. Zainstaluj do środowiska niewłaściwego projektu, a czysty
checkout kolegi nie uruchomi twojego kodu, bez żadnego komunikatu, który
nazwałby przyczynę. Reguła, która usuwa całą tę klasę: \textbf{nigdy nie
uruchamiaj \code{pip install} w tej książce.} Zależności trafiają do
\code{pyproject.toml} przez \code{uv add}, a środowiska buduje się z pliku
blokady.
\end{trapbox}

\mermaidfig{env-is-a-directory}{Gdzie oba ekosystemy kładą zależność i co z tego wynika}{katalog, nie maszyna}

\section{Dwa pliki, a prawdą jest tylko jeden}
\label{sec:ch02-two-files}
\index{pakowanie!pyproject.toml}
\index{pakowanie!plik blokady}

\noindent
\code{pyproject.toml} to plik projektu, a ten należący do tej książki jest
szkieletem obiecanym w opisie rozdziału. Oto część mówiąca, czego kod
potrzebuje, żeby działać:

\pyregion{code/pyproject.toml}{project}{Tabela \code{[project]}: czego ten projekt potrzebuje, żeby działać}{lst:ch02-project}

\noindent
I, osobno, czego ty potrzebujesz, żeby nad nim pracować:

\pyregion{code/pyproject.toml}{groups}{Grupy zależności: czego potrzebujesz do pracy, a to nie ta sama lista}{lst:ch02-groups}

\noindent
Ten podział to pierwsza rzecz do podkradnięcia. \code{[project].dependencies}
to to, co instaluje wdrożenie; grupa zależności nie jest instalowana, dopóki
ktoś o nią nie poprosi. Obraz kontenera zbudowany z pierwszej listy nie
zawiera runnera testów, lintera ani checkera typów, i nie zawiera ich dlatego,
gdzie napisano linię, a nie dlatego, że ktoś pamiętał, żeby je potem wyciąć.

\begin{csbox}
\code{pyproject.toml} to \code{.csproj}: tożsamość projektu, jego zależności i
konfiguracja każdego narzędzia, które na nim działa. Jest w TOML-u zamiast w
XML-u i mieści znacznie więcej \dash{} \pkg{ruff}, \pkg{pytest}, \pkg{pyright}
i \pkg{mypy} konfiguruje się w tym jednym pliku, więc repozytorium Pythona nie
ma odpowiednika rozproszenia na \code{.editorconfig},
\code{Directory.Build.props} i zestaw reguł.
\end{csbox}

\noindent
Teraz drugie pytanie, i znowu odpowiedz przed czytaniem dalej. Ten projekt
deklaruje zależności, które właśnie widziałeś, w dwóch krótkich listach. Gdy
\code{uv sync} skończy, ile dystrybucji będzie w \code{.venv}?

\pyfile{code/ch02/lockfile.py}{Liczenie tego, co zadeklarowano, przeciw temu, co rozwiązano}{lst:ch02-lockfile}

\transcript{ch02-lockfile}

\noindent
Większości środowiska nikt nigdy nie nazwał. To są zależności przechodnie
\dash{} czego potrzebuje \pkg{fastapi}, czego potrzebuje \pkg{anthropic},
czego z kolei potrzebują tamte \dash{} a plikiem, który przypina je do
dokładnych wersji, jest \code{uv.lock}, i on jest commitowany. To jest
rozróżnienie, na którym stoi reszta rozdziału: \code{pyproject.toml} zapisuje
\emph{zamiar}, a \code{uv.lock} zapisuje \emph{wynik}. Tylko jednego z nich da
się użyć do odtworzenia środowiska.

\begin{csbox}
\code{uv.lock} to \code{packages.lock.json}, a argument za commitowaniem go
już akceptujesz. Różnica jest raczej kulturowa niż techniczna: blokowanie jest
w .NET opcjonalne i rzadkie, a w Pythonie \code{uv.lock} powstaje domyślnie
przy pierwszym uruchomieniu czegokolwiek, więc jedynym sposobem, żeby go nie
mieć, jest celowe dopisanie go do \code{.gitignore}.
\end{csbox}

\begin{trapbox}
\textbf{\code{requirements.txt} to plik blokady.} To lista życzeń, i zwykle
niepełna. To plik, który napisał jakiś człowiek, zwykle wymieniając tylko
pakiety, o których pomyślał, zwykle z zakresami zamiast dokładnych wersji, i
nic nie sprawdza go względem tego, co jest zainstalowane. Dwie maszyny idące
za tym samym \code{requirements.txt} rutynowo kończą z różnymi wersjami
większości tego, co uruchamiają, a awaria pojawia się tygodnie później jako
błąd odtwarzalny na jednej z nich. \code{uv.lock} jest generowany, obejmuje
cały graf, zapisuje skrót każdego artefaktu, a \code{uv sync --locked} odmawia
instalacji czegokolwiek innego.
\end{trapbox}

\mermaidfig{declare-resolve-lock}{Zadeklarowane, rozwiązane, zainstalowane: trzy pliki i ten, który odtwarza środowisko}{deklaracja, rozwiązanie, blokada}

\begin{exercise}{e02_01_pinned_version}{Czytaj blokadę, nie listę życzeń}
Plik startowy ma jedną funkcję do dokończenia: mając sparsowany
\code{uv.lock} i nazwę dystrybucji, zwróć dokładną wersję, do której jest
przypięta. Testy napędzają ją plikiem blokady tej książki, więc jeden z nich
prosi o pakiet, który nie występuje w żadnej liście zależności w całym
repozytorium \dash{} i mimo to oczekuje, że znajdziesz dla niego wersję. Ten
test jest całym rozdziałem w jednej asercji.
\end{exercise}

\section{Uruchamiaj, nie aktywuj}
\label{sec:ch02-run}
\index{pakowanie!uv run}

\noindent
Każdy tutorial, jaki przeczytasz, każe ci aktywować środowisko. Odruch, który
sprawia, że inżynierowi .NET jest z tym wygodnie, warto przesłuchać, zanim go
przyjmiesz: które z tego, co jutro uruchomi twój kod, wykona ten krok
aktywacji?

Powłoka, którą otworzyłeś, i nic poza tym. Aktywacja edytuje \code{PATH} w
jednym procesie. Drugi terminal jej nie wykonał. Cel w \code{Makefile} jej nie
wykonał. \code{cron} jej nie wykonał. CI jej nie wykonało, ani runner testów w
twoim edytorze, dlatego właśnie zestaw testów, który przechodzi w terminalu, a
pada w IDE, to tak często to i nic ciekawszego.

Cały przepływ pracy tej książki to trzy polecenia, z których żadne nie
potrzebuje aktywnego środowiska:

\begin{shellcmd}
uv init --name incident-tools  # a project, with a src/ layout
uv add httpx                   # a dependency, written into pyproject.toml
uv run incident-tools          # and run its entry point, in that environment
\end{shellcmd}

\noindent
\code{uv run} czyta \code{pyproject.toml} i \code{uv.lock}, upewnia się, że
\code{.venv} im odpowiada, a potem uruchamia twoje polecenie w środku. Robi to
przy każdym wywołaniu, dlatego jest poprawne z powłoki, z \code{Makefile}, z
\code{cron}-a i z CI, bez niczego do zapamiętania i bez niczego do
ustawienia. \code{uvx} to jego rodzeństwo dla narzędzia, które chcesz
uruchomić, a nie mieć w zależnościach \dash{} odpowiednik \code{dotnet tool},
bez kroku instalacji.

\begin{csbox}
\code{uv run} to \code{dotnet run}, i przywraca zależności przed
uruchomieniem z tego samego powodu. \code{uv add} to \code{dotnet add
package}. Tym bez odpowiednika jest aktywacja, a to dlatego, że .NET nigdy jej
nie potrzebował: nie ma czego aktywować, gdy środowiskiem jest maszyna.
\end{csbox}

\begin{trapbox}
\textbf{Aktywuję środowisko i dopiero potem uruchamiam.} Ten nawyk nie tyle
jest błędny, co zawodny, i zawodzi w kierunku najtrudniejszym do
zdiagnozowania: działa na twojej maszynie. Aktywacja jest per powłoka, więc
awarie lądują tam, gdzie nie możesz patrzeć \dash{} w zadaniu z
harmonogramu, na agencie buildu, przy pierwszym checkoucie kolegi. Zamiast
tego dopisz \code{uv run} przed poleceniem, a pytanie, w którym środowisku
jesteś, przestanie istnieć.
\end{trapbox}

\begin{warning}
\code{uv run} i \code{uv sync --locked} robią co innego z plikiem blokady,
który został w tyle za \code{pyproject.toml}, i ta różnica jest warta
poznania. Dopisz zależność ręcznie, a potem uruchom \code{uv sync --locked}:
kończy kodem niezerowym i mówi ci, że blokada wymaga aktualizacji. Uruchom
zamiast tego \code{uv run}, a przeblokuje, zainstaluje i pójdzie dalej bez
słowa. To jest wygodne przy biurku i błędne w CI, gdzie cały sens polega na
tym, żeby zawieść, a nie po cichu zainstalować coś, czego nikt nie
przejrzał. W każdym kontekście zautomatyzowanym używaj
\code{uv sync --locked}.
\end{warning}

\mermaidfig{run-not-activate}{Aktywacja jest per powłoka; \code{uv run} jest per wywołanie}{uruchamiaj, nie aktywuj}

\section{Co daje uv, zmierzone}
\label{sec:ch02-measured}
\index{pomiar!E2}

\noindent
To, że \pkg{uv} jest szybkie, jest pierwszą rzeczą, którą każdy o nim mówi, i
samo w sobie nie jest powodem, żeby przyjąć narzędzie. Pytaniem wartym
odpowiedzi jest \emph{gdzie idzie czas}, bo to ono rozstrzyga, czy
skonfigurowanie cache'u w CI jest warte popołudnia.

Zatem: własny plik blokady tej książki, wyeksportowany tak, by oba narzędzia
instalowały dokładnie te same \val{e2.packages} dystrybucji, razem ze
skrótami, bez niczego do rozwiązania dla żadnego z nich, na przypiętym
interpreterze, każde z prywatnym cache'em pakietów
opróżnianym przed próbą zimną i zachowywanym przed ciepłą. Środowisko
powstaje w mierzonym odcinku w obu przypadkach, bo przy \code{pip} oznacza to
\code{python -m venv}, a czytelnik kupuje całą drogę od checkoutu do gotowego
do uruchomienia. Mediana z \val{e2.trials} prób.

\begin{center}
\small
\begin{tabularx}{\linewidth}{@{}Xrrr@{}}
\toprule
& \textbf{cache zimny} & \textbf{cache ciepły} & \textbf{co daje cache} \\
\midrule
\code{uv sync --locked} & \val{e2.uv.cold}\,s & \val{e2.uv.warm}\,s & \val{e2.uv.cachegain}$\times$ \\
\code{pip install -r} & \val{e2.pip.cold}\,s & \val{e2.pip.warm}\,s & \val{e2.pip.cachegain}$\times$ \\
\midrule
\textbf{stosunek} & \val{e2.cold.ratio}$\times$ & \val{e2.warm.ratio}$\times$ & \\
\bottomrule
\end{tabularx}
\end{center}

\noindent
Ostatnia kolumna jest wynikiem, i nie tym, który sugeruje nagłówek. Zapełniony
cache daje \pkg{uv} czynnik \val{e2.uv.cachegain}; \code{pip}-owi nie daje
prawie nic. Czas \code{pip}-a to nie pobieranie. To własna praca \code{pip}-a
\dash{} budowanie środowiska, rozpakowywanie kół, kopiowanie plików \dash{} i
żaden cache nie usuwa z niej niczego, dlatego ciepły przebieg \code{pip}-a
jest szybszy od zimnego mniej więcej o jedną dziesiątą, a nie o rząd
wielkości. Te dwa cache nie trzymają tego samego.
\code{pip} trzyma każde koło wciąż spakowane, więc ciepła instalacja
rozpakowuje je wszystkie od nowa; \pkg{uv} trzyma pliki już rozpakowane i
materializuje z nich środowisko bez kopiowania bajtów, co nazywa
\code{--link-mode} i co na Linuksie domyślnie jest klonem
kopiuj-przy-zapisie. Gdy koła są już lokalne, \pkg{uv} nie ma prawie nic do
zrobienia.

Dwie konsekwencje, na które możesz zadziałać. Cache'owanie katalogu pakietów w
CI jest warte zachodu przy \pkg{uv} i niemal bezcelowe przy \code{pip}. A
build kontenera instalujący \code{pip}-em płaci ten sam koszt przy każdym
przebudowaniu warstwy, choćbyś cache'ował jak najlepiej.

\begin{note}
\textbf{Czym te liczby są, a czym nie są.} Stosunki są połową przenośną;
sekundy to jedna maszyna jednego popołudnia, a maszyna, która je wykonała, ma
wolną i zawodną sieć. Trzy zimne próby \code{pip}-a wylądowały między
\val{e2.pip.cold.best} a \val{e2.pip.cold.worst} sekundy, co jest ciasne.
Najlepsza zimna próba \pkg{uv} to \val{e2.uv.cold.best} sekundy, a najgorsza
\val{e2.uv.cold.worst} sekund \dash{} pobieranie, które zatrzymało się na ponad pięć
minut. To jest właściwość łącza, a nie narzędzia, dlatego raportowana jest
mediana zamiast średniej, i dlatego liczba zimna jest najmniej przenośną
pozycją w tabeli. Skrypt to \code{code/measure/e2\_install.py}, a surowe próby
są zacommitowane obok niego; uruchom go sam i spodziewaj się innych sekund i
podobnych stosunków. \code{pip} był w wersji \valtext{e2.pip.version}, tej,
która jedzie z przypiętym interpreterem.
\end{note}

\section{Szkielet do skopiowania}
\label{sec:ch02-skeleton}
\index{pakowanie!układ src}
\index{pakowanie!ruff}

\noindent
\code{uv init} daje ci projekt, którego pakiet siedzi w katalogu \code{src/}, a
nie w korzeniu repozytorium, i powód wart jest akapitu, bo wygląda na
ceremonię, a nią nie jest. Jeśli twój pakiet leży w korzeniu, da się go
zaimportować z korzenia niezależnie od tego, czy został zainstalowany, więc
twoje testy przechodzą względem plików na dysku i nie mówią nic o tym, czy
rzecz, którą wysyłasz, by działała. Pod \code{src/} musi być zainstalowany,
żeby dało się go zaimportować, co znaczy, że testy idą względem
zainstalowanego artefaktu. To różnica między testowaniem drzewa źródeł a
testowaniem pakietu.

\code{uv init} zapisuje też punkt wejścia: nazwę w \code{[project.scripts]}
związaną z funkcją, która staje się poleceniem na ścieżce, gdy projekt zostanie
zainstalowany. To jest \code{Program.cs} projektu Pythona, tyle że projekt
może zadeklarować ich kilka i każdy jest jedną linią.

\begin{csbox}
W Pythonie nie ma \code{Program.cs} ani \code{Main}. Moduł, który ma dać się
uruchomić, kończy się na \code{if \_\_name\_\_ == "\_\_main\_\_":}, co jest
tematem Rozdziału~\ref{ch:runtime}; pakiet, który ma zainstalować polecenie,
deklaruje je w \code{[project.scripts]}. Te dwie rzeczy są niezależne, a
większość projektów chce obu.
\end{csbox}

\noindent
Ostatnim elementem jest linter. \pkg{ruff} zastępuje naraz analizatory i
\code{dotnet format} \dash{} jedno narzędzie, konfigurowane w
\code{pyproject.toml}, które i raportuje, i poprawia \dash{} i jest na tyle
szybkie, że nie ma powodu, żeby nie uruchamiać go przy każdym zapisie.
\pkg{pre-commit} jest tym, co czyni je nieuniknionym dla wszystkich
pozostałych: hak gita, zadeklarowany w pliku, który puszcza formatter i linter
po plikach dotkniętych commitem i odmawia commita, gdy są niezadowolone. Nic w
.NET nie robi tego rutynowo, bo build już to robi.

Własny szkielet tej książki puszcza linter i checker w CI, a nie w haku, i to
rozróżnienie warto mieć: hak żyje na maszynie jednego programisty i da się go
pominąć flagą, a buildu nie. Uczciwy układ to jedno i drugie, przy czym hak
jest tą połową, która oszczędza ci rundy, a nie tą, na której można polegać.

\begin{exercise}{e02_02_install_set}{Co wdrożenie naprawdę instaluje?}
Dokończ funkcję, która odpowiada na pytanie, dla którego istnieje podział na
\code{[project]} i \code{[dependency-groups]}: mając sparsowany
\code{pyproject.toml} i wybór grup, które dystrybucje zostałyby
zainstalowane? Jeden test napędza ją małym dokumentem niosącym każdy kształt,
jaki może przyjąć wymaganie; ostatni pyta własny projekt tej książki, czy jego
wdrożenie zawierałoby \pkg{pytest}.
\end{exercise}

\noindent
Masz teraz środowisko, któremu możesz zaufać, i sposób na odbudowanie go z
pliku. To jest podłoga, na której stoją listingi każdego kolejnego rozdziału.
Następny jest o innym rodzaju zaufania: Rozdział~\ref{ch:typing} jest tam,
gdzie twój odruch, że adnotacja typu coś znaczy, spotyka środowisko
uruchomieniowe, które nigdy żadnej nie sprawdza.
